% GENERATED FILE. Edit canonical lessons, then run npm run generate:books.
% canonical-source-hash: fnv1a64:6fb7bd94a12b36d4
% canonical-lessons: BN-C273-doinik, BN-W273-letter-oi-kar, BN-C273-cheye, BN-C273-dhore, BN-R273-comparing-and-time

\chapter{Than, Most, and How Long}
\label{ch:bn-than-most-and-how-long}
\hlchaptermodality{273}

\begin{chapteropening}
\textbf{By the end of this chapter:} \emph{I can compare (-er cheye, sôbcheye), say how long something has gone on (dui bôchhôr dhôre, sôkāl theke rāt pôrjôntô), and read and write the sign \bn{ৈ} in \bn{দৈনিক}.}

Complete BN-R273-comparing-and-time: compare things and say how long.
\end{chapteropening}

\section[\texorpdfstring{\bn{দৈনিক} (doinik)}{দৈনিক (doinik)}]{\bn{দৈনিক} (doinik) — of each day; a daily paper}

\label{lesson:BN-C273-doinik}



\begin{quote}
Before the new one: say \textquotedblleft{}I want to be a doctor\textquotedblright{}. (\emph{āmi ḍāktār hôte chāi}.) Then \textbf{\bn{রোজ}}, \textquotedblleft{}daily, every day\textquotedblright{}.
\end{quote}

\subsection*{You'll want to know: \bn{দৈনিক}}
\textbf{\bn{দৈনিক}} — \emph{doinik} — \textquotedblleft{}of each day; a daily paper\textquotedblright{}.

\begin{quote}
\textbf{\bn{আমার} \bn{দৈনিক} \bn{খরচ} \bn{একশো} \bn{টাকা।}}
\end{quote}

\begin{quote}
\emph{āmār doinik khôrôch eksho ṭākā.} — \textquotedblleft{}My daily spending is a hundred taka.\textquotedblright{}
\end{quote}

\textbf{\bn{রোজ}} tells you when: \textquotedblleft{}every day\textquotedblright{}. \textbf{\bn{দৈনিক}} describes a thing: \textbf{\bn{দৈনিক} \bn{খরচ}}, \textquotedblleft{}daily spending\textquotedblright{}, \textbf{\bn{দৈনিক} \bn{কাজ}}, \textquotedblleft{}daily work\textquotedblright{}. On its own it is a newspaper that comes out every day: \textbf{\bn{বাবা} \bn{রোজ} \bn{একটা} \bn{দৈনিক} \bn{পড়েন}} (\emph{bābā roj ekṭā doinik pôṛen}), \textquotedblleft{}Father reads a daily every day\textquotedblright{}.

\begin{cousinweb}
\textbf{\bn{দৈনিক}} (\emph{doinik}) is a tatsama: Sanskrit \emph{dainika}, \textquotedblleft{}daily\textquotedblright{} (Monier-Williams). It belongs with \emph{dina}, \textquotedblleft{}day\textquotedblright{}, Bengali \textbf{\bn{দিন}}: CDIAL (6566) gives the related \emph{daina}, \textquotedblleft{}relating to a day\textquotedblright{}, and refers it to \emph{dina}; \emph{dainika} has the same strengthened first vowel, \emph{ai}, with \emph{-ika}, \textquotedblleft{}of, belonging to\textquotedblright{}, added.
\end{cousinweb}

\subsection*{Your turn}
Say these aloud:

\begin{itemize}
\raggedright
  \item \emph{doinik}
  \item \textquotedblleft{}My daily spending is a hundred taka\textquotedblright{} — \emph{āmār doinik khôrôch eksho ṭākā}
  \item \textquotedblleft{}Father reads a daily every day\textquotedblright{} — \emph{bābā roj ekṭā doinik pôṛen}
\end{itemize}

\begin{tcolorbox}[breakable,colback=teal!4,colframe=teal!35!black,title={Before you move on}]
What does \textbf{\bn{দৈনিক}} describe? (A thing of each day.) Which word says \textquotedblleft{}every day\textquotedblright{}, when? (\textbf{\bn{রোজ}}.)
\end{tcolorbox}
\section[\texorpdfstring{\bn{ৈ} (-oi)}{ৈ (-oi)}]{\bn{ৈ} (-oi) — the sign in front of the \bn{দ} of \bn{দৈনিক}}

\label{lesson:BN-W273-letter-oi-kar}



\begin{quote}
Before the new one: say \textbf{\bn{দৈনিক}}. Then write \textbf{\bn{দে}} (\emph{de}) and \textbf{\bn{নৌ}} (\emph{nou}).
\end{quote}

\subsection*{Script you'll notice: \bn{ৈ}}
\textbf{\bn{ৈ}} — \emph{-oi}, the vowel of \textbf{\bn{দৈনিক}}. Like \textbf{\bn{ে}}, it is written in front of its consonant, though it is said after it.

Set the print of \textbf{\bn{দৈ}} beside \textbf{\bn{দে}}. \textbf{\bn{দে}} has \textbf{\bn{ে}} in front: an upright curve with a small knob at its foot. \textbf{\bn{দৈ}} has the same, and above it a hook stands on the headline, bent over to the left, its tip at the upper left: the same hook as the one on the stroke of \textbf{\bn{ৌ}}.

\subsection*{Writing: \bn{ৈ} — copy the printed shape}
Put your pen beside \textbf{\bn{দৈ}} and copy the shape you can see, slowly and large: the \textbf{\bn{ে}} with its hook above the headline, then the \textbf{\bn{দ}}. Check the copy against \textbf{\bn{দে}}: only the hook differs.

\begin{quote}
This book does not yet tell you where to start \bn{ৈ} or which way the pen goes. That is not written down here until it can be written down with a source. Copying the printed shape needs no such source, and it is how the shape becomes yours in the meantime.
\end{quote}

\subsection*{Your turn}
\begin{itemize}
\raggedright
  \item \emph{Find:} \bn{ৈ} in front of the \bn{দ} of \textbf{\bn{দৈনিক}}
  \item \emph{Write it:} \textbf{\bn{দে}}, then \textbf{\bn{দৈ}}, from the print
  \item \emph{Write it:} \textbf{\bn{দৈনিক}}
  \item \emph{Write it:} \textbf{\bn{অভিজ্ঞতা}}, from memory
\end{itemize}

\begin{tcolorbox}[breakable,colback=teal!4,colframe=teal!35!black,title={Before you move on}]
Which sign is in front of the \textbf{\bn{দ}} of \textbf{\bn{দৈনিক}}? (\textbf{\bn{ৈ}}.) How is it unlike \textbf{\bn{ে}}? (A hook stands above the headline.) Which other sign has that hook? (\textbf{\bn{ৌ}}, on its stroke after the consonant.)
\end{tcolorbox}
\section[\texorpdfstring{-\bn{এর} \bn{চেয়ে} / \bn{সবচেয়ে} (-er cheye / sôbcheye)}{-এর চেয়ে / সবচেয়ে (-er cheye / sôbcheye)}]{-\bn{এর} \bn{চেয়ে} / \bn{সবচেয়ে} (-er cheye / sôbcheye) — than, and most}

\label{lesson:BN-C273-cheye}



\begin{quote}
Say \textquotedblleft{}my daily spending\textquotedblright{}. (\emph{āmār doinik khôrôch}.) Then \textbf{\bn{বড়}} and \textbf{\bn{ছোট}}, \textquotedblleft{}big\textquotedblright{} and \textquotedblleft{}small\textquotedblright{}.
\end{quote}

\subsection*{You'll want to know: than}
\begin{quote}
\textbf{\bn{নদী} \bn{পুকুরের} \bn{চেয়ে} \bn{বড়।}}
\end{quote}

\begin{quote}
\emph{nôdī pukurer cheye bôṛo.} — \textquotedblleft{}A river is bigger than a pond.\textquotedblright{}
\end{quote}

The thing you measure against takes the \textbf{-\bn{র}} form and \textbf{\bn{চেয়ে}} (\emph{cheye}), \textquotedblleft{}than\textquotedblright{}. The describing word does not change: \textbf{\bn{বড়}} is \textquotedblleft{}big\textquotedblright{} and \textquotedblleft{}bigger\textquotedblright{}. \textbf{\bn{ট্রেন} \bn{বাসের} \bn{চেয়ে} \bn{তাড়াতাড়ি} \bn{যায়}} (\emph{ṭren bāser cheye tāṛātāṛi jāy}), \textquotedblleft{}a train goes faster than a bus\textquotedblright{}.

\subsection*{You'll want to know: most}
\textbf{\bn{সব}} (\emph{sôb}) is \textquotedblleft{}all\textquotedblright{}, so \textbf{\bn{সবচেয়ে}} (\emph{sôbcheye}) is \textquotedblleft{}than all\textquotedblright{}: \textquotedblleft{}most, the -est\textquotedblright{}.

\begin{quote}
\textbf{\bn{আমার} \bn{সবচেয়ে} \bn{প্রিয়} \bn{খাবার} \bn{মাছ।}}
\end{quote}

\begin{quote}
\emph{āmār sôbcheye priyô khābār māchh.} — \textquotedblleft{}My favourite food of all is fish.\textquotedblright{}
\end{quote}

\textbf{\bn{এটা} \bn{সবচেয়ে} \bn{দামি} \bn{বই}} (\emph{eṭā sôbcheye dāmi bôi}), \textquotedblleft{}this is the most expensive book\textquotedblright{}.

\subsection*{Your turn}
Say these aloud:

\begin{itemize}
\raggedright
  \item \textquotedblleft{}The city is bigger than the village\textquotedblright{} — \emph{shôhôr grāmer cheye bôṛo}
  \item \textquotedblleft{}A train goes faster than a bus\textquotedblright{} — \emph{ṭren bāser cheye tāṛātāṛi jāy}
  \item \textquotedblleft{}This is the most expensive book\textquotedblright{} — \emph{eṭā sôbcheye dāmi bôi}
\end{itemize}

\begin{tcolorbox}[breakable,colback=teal!4,colframe=teal!35!black,title={Before you move on}]
\textquotedblleft{}Than\textquotedblright{}? (\textbf{-\bn{এর} \bn{চেয়ে}}.) Does \textbf{\bn{বড়}} change for \textquotedblleft{}bigger\textquotedblright{}? (No.) \textquotedblleft{}Most\textquotedblright{}? (\textbf{\bn{সবচেয়ে}}: \textquotedblleft{}than all\textquotedblright{}.)
\end{tcolorbox}
\section[\texorpdfstring{… \bn{ধরে} / \bn{থেকে} … … (… dhôre / theke …)}{… ধরে / থেকে … … (… dhôre / theke …)}]{… \bn{ধরে} / \bn{থেকে} … \bn{পর্যন্ত} (dhôre / theke … pôrjôntô) — how long}

\label{lesson:BN-C273-dhore}



\begin{quote}
Say \textquotedblleft{}A river is bigger than a pond\textquotedblright{}. (\emph{nôdī pukurer cheye bôṛo}.) Then \textbf{\bn{ধরা}}, \textquotedblleft{}to hold\textquotedblright{}.
\end{quote}

\subsection*{You'll want to know: for a stretch of time}
\begin{quote}
\textbf{\bn{আমি} \bn{দুই} \bn{বছর} \bn{ধরে} \bn{বাংলা} \bn{শিখছি।}}
\end{quote}

\begin{quote}
\emph{āmi dui bôchhôr dhôre bānglā shikhchhi.} — \textquotedblleft{}I have been learning Bengali for two years.\textquotedblright{}
\end{quote}

\textbf{\bn{ধরে}} is \textbf{\bn{ধরা}} in the \emph{-e} form: \textquotedblleft{}holding\textquotedblright{} a stretch of time. With the going-on present, it says what started then and is still going on: English says \textquotedblleft{}have been\textquotedblright{}, Bengali says \textquotedblleft{}am\textquotedblright{}. In the past it is a stretch that is over: \textbf{\bn{আমরা} \bn{তিন} \bn{দিন} \bn{ধরে} \bn{হাঁটলাম}} (\emph{āmrā tin din dhôre hā̃ṭlām}), \textquotedblleft{}we walked for three days\textquotedblright{}.

\subsection*{You'll want to know: from … until}
\textbf{\bn{থেকে}}, \textquotedblleft{}from\textquotedblright{}, which you have met, is \textbf{\bn{থাকা}} in the \emph{-e} form; \textbf{\bn{পর্যন্ত}} (\emph{pôrjôntô}) is \textquotedblleft{}until, as far as\textquotedblright{}. They frame a stretch with two ends:

\begin{quote}
\textbf{\bn{দোকান} \bn{সকাল} \bn{থেকে} \bn{রাত} \bn{পর্যন্ত} \bn{খোলা} \bn{থাকে।}}
\end{quote}

\begin{quote}
\emph{dokān sôkāl theke rāt pôrjôntô kholā thāke.} — \textquotedblleft{}The shop is open from morning until night.\textquotedblright{}
\end{quote}

\textbf{\bn{পর্যন্ত}} measures places too: \textbf{\bn{বাড়ি} \bn{থেকে} \bn{স্টেশন} \bn{পর্যন্ত}} (\emph{bāṛi theke sṭeshôn pôrjôntô}), \textquotedblleft{}from the house as far as the station\textquotedblright{}.

\subsection*{Your turn}
Say these aloud:

\begin{itemize}
\raggedright
  \item \textquotedblleft{}I have been learning Bengali for two years\textquotedblright{} — \emph{āmi dui bôchhôr dhôre bānglā shikhchhi}
  \item \textquotedblleft{}We walked for three days\textquotedblright{} — \emph{āmrā tin din dhôre hā̃ṭlām}
  \item \textquotedblleft{}I work from morning until evening\textquotedblright{} — \emph{āmi sôkāl theke sôndhyā pôrjôntô kāj kôri}
\end{itemize}

\begin{tcolorbox}[breakable,colback=teal!4,colframe=teal!35!black,title={Before you move on}]
\textquotedblleft{}For two years\textquotedblright{}? (\textbf{\bn{দুই} \bn{বছর} \bn{ধরে}}.) \textquotedblleft{}Have been learning\textquotedblright{}? (The going-on present: \textbf{\bn{শিখছি}}.) \textquotedblleft{}From … until\textquotedblright{}? (\textbf{\bn{থেকে} … \bn{পর্যন্ত}}.)
\end{tcolorbox}
\section[(review)]{Review — than, most, and how long}

\label{lesson:BN-R273-comparing-and-time}



\begin{quote}
Say \textquotedblleft{}The city is bigger than the village\textquotedblright{}, then \textquotedblleft{}for two years\textquotedblright{}. (\emph{shôhôr grāmer cheye bôṛo}; \emph{dui bôchhôr dhôre}.)
\end{quote}

\subsection*{Your turn}
\begin{itemize}
\raggedright
  \item \emph{Say it:} The mountains are more beautiful than the city. I'd like to go there.
  \item \emph{Say it:} This is my favourite book of all.
  \item \emph{Say it:} Mira has been working at the hospital for three years.
  \item \emph{Say it:} The market is open from morning until evening.
  \item \emph{Say it:} Father reads a daily every morning.
  \item \emph{Write it:} \textbf{\bn{দৈনিক}}, then \textbf{\bn{নৌকা}}
\end{itemize}

\begin{tcolorbox}[breakable,colback=teal!4,colframe=teal!35!black,title={Before you move on}]
\textquotedblleft{}Than\textquotedblright{}? (\textbf{-\bn{এর} \bn{চেয়ে}}.) \textquotedblleft{}Most\textquotedblright{}? (\textbf{\bn{সবচেয়ে}}.) \textquotedblleft{}For three years\textquotedblright{}? (\textbf{\bn{তিন} \bn{বছর} \bn{ধরে}}.) Which sign is in front of the \textbf{\bn{দ}} of \textbf{\bn{দৈনিক}}? (\textbf{\bn{ৈ}}.)
\end{tcolorbox}
